\documentclass[12pt,a4paper]{ctexart}
\usepackage[top=2.5cm,bottom=2.5cm,left=2.5cm,right=2.5cm]{geometry}
\usepackage{amsmath,amssymb,amsthm,physics}
\usepackage{graphicx,float,booktabs,array,caption,multirow}
\usepackage{hyperref,xcolor,enumitem,mathtools}
\usepackage[numbers,sort&compress]{natbib}
\hypersetup{colorlinks=true,linkcolor=blue,citecolor=blue,urlcolor=blue}
\theoremstyle{definition}
\newtheorem{definition}{定义}[section]
\newtheorem{theorem}{定理}[section]
\newtheorem{remark}{注记}[section]

\title{\heiti 格点QCD中的Wilson线：\\
从禁闭势到部分子物理的非定域探针}
\author{基于本库全部相关文献与教材的综合解析}
\date{\today}

\begin{document}
\maketitle

\begin{abstract}
\noindent
Wilson线（Wilson line）——即规范场沿路径的编序指数积分——是规范场论中最基本的非定域客体。在格点QCD中，Wilson线以规范链接（gauge link）变量的形式成为构建整个理论的基础砖块：它是保持格点费米子作用规范不变性的平行传输子，是定义Wilson圈并从中提取静态夸克-反夸克势的核心要素，是通过Polyakov圈探测禁闭-退禁闭相变的有序参量，更是现代部分子物理中准PDF和准TMD-PDF算符的不可分割的组成部分。

本文基于本库中的三本核心教材（Gattringer与Lang的\textit{Quantum Chromodynamics on the Lattice}、Gupta的\textit{Introduction to Lattice QCD}、Peskin与Schroeder的\textit{An Introduction to Quantum Field Theory}）以及约35篇相关研究论文，从四个层次系统解析格点QCD中的Wilson线：

\textbf{第一层——基本定义与几何：}从连续规范场中的路径编序指数$U_\Gamma = \mathcal{P}\exp(ig\int_\Gamma dx^\mu A_\mu)$到格点上的规范链接变量$U_\mu(n)$，阐明Wilson线作为``平行传输子''的本质——将规范群在一点的颜色参考系平行传输到另一点。讨论了格点上规范链接的规范变换性质和由之构造的两种基本规范不变量：闭合Wilson圈和连接费米子-反费米子的开Wilson线。

\textbf{第二层——静态夸克势：}详述Wilson圈（由两条空间Wilson线和两条时间传输子构成的矩形闭合回路）如何通过时间关联函数的量子力学诠释，给出静态夸克-反夸克势$V(r) = A + B/r + \sigma r$。追溯了强耦合展开中面积律（area law）$\langle W_C\rangle \propto \exp(-\sigma \cdot \text{Area})$的推导——这是格点QCD证明线性禁闭势的经典结果。在此基础上引入了Sommer参数$r_0\simeq 0.5$~fm作为格距标定的标准方法。同时讨论了Polyakov圈（热力学Wilson线）作为有限温度下禁闭-退禁闭相变有序参量的物理诠释。

\textbf{第三层——部分子物理中的Wilson线：}系统阐述LaMET框架下准PDF和准TMD-PDF算符中Wilson线的构造。对于准PDF，算符$\bar\psi(z\hat n_z)\Gamma U_z(z,0)\psi(0)$中的直Wilson线$U_z$编码了从$0$到$z$的规范平行传输。对于准TMD-PDF，staple型（订书钉形）Wilson线$W_{\sqsubset}(b_\perp, L, z)$需要三段Wilson线的组合来同时捕获纵向和横向的规范平移，其$L\to\infty$极限模拟了光锥方向上的无穷长Wilson线。

\textbf{第四层——线性发散与重整化：}这是现代格点部分子物理中最关键的技术问题。Wilson线的自能修正产生与Wilson线长度成正比的线性发散$e^{-\delta\bar{m}L/a}$，在连续极限下发散。本文详细分析了三种应对策略：（1）质量抵消项形式——利用辅助$z$场形式，线性发散等价于重夸克有效理论中的质量重整化$\delta m$；（2）Wilson圈减除——利用矩形Wilson圈平方根$\sqrt{Z_E}$消除准TMD-PDF中的线性发散、pinch-pole奇点和cusp发散；（3）自重整化方法——通过同一矩阵元在不同$z$处比值消除线性发散而不依赖微扰计算。对比了RI/MOM方案的失败及其根本原因——胶子-规范链接顶点处的非微扰线性发散在离壳减除中未被完全消除。

最后，本文讨论了HYP涂抹（Hypercubic Smeared Link）对Wilson线的改进效果、规范固定精度对长Wilson线的影响（定量给出$\delta_F = 10^{-8}$精度下在$a=0.03$~fm格距上1.2~fm Wilson线可达12\%偏差的实证数据），以及Wilson线在重夸克有效理论（HQET）中的核心角色。

全篇以统一的数学语言（路径编序指数、规范变换性质、谱表示）将分散在本库各文献中的Wilson线知识编织成一个有机整体，为读者提供了从经典禁闭物理到前沿部分子结构的完整Wilson线理论图景。
\end{abstract}

\tableofcontents
\newpage

%==========================================================================
\section{引言：Wilson线——规范理论的脊梁}
%==========================================================================

\subsection{连续规范场中的路径编序指数}

在连续的非Abel规范场论中，Wilson线定义为规范场沿时空路径$\Gamma$的\textbf{路径编序（path-ordered）指数积分}：

\begin{definition}[连续Wilson线]
设$\Gamma$为从$x$到$y$的一条有向路径，参数化为$z^\mu(s)$，$s\in[0,1]$，$z(0)=x$，$z(1)=y$。路径编序指数定义为：
\begin{equation}
\boxed{U_\Gamma(y,x) \equiv \mathcal{P}\exp\!\left[ig\int_0^1 ds\, \frac{dz^\mu(s)}{ds} A_\mu(z(s))\right]},
\label{eq:wilson_line_continuum}
\end{equation}
其中$A_\mu(x) = A_\mu^a(x)T^a$是李代数取值的规范势（$T^a$是规范群的生成元），$\mathcal{P}$表示路径编序算符——沿路径$\Gamma$的方向，越早出现的因子越靠\textbf{右}。
\end{definition}

\textbf{路径编序的显式定义：}将路径划分为$N$个无穷小线段$\Delta z_k$，则
\begin{equation}
\mathcal{P}\exp\!\left[ig\int_\Gamma dx^\mu A_\mu\right] \equiv \lim_{N\to\infty} \prod_{k=1}^N \exp\!\big[ig\,\Delta z_k^\mu A_\mu(z_k)\big],
\label{eq:path_ordering}
\end{equation}
其中乘积按$k$增大方向从左到右排列。由于$A_\mu$的不同分量在非Abel理论中不对易，路径编序是\textbf{必需的}——没有它，指数积分将失去规范的协变性。

\subsection{规范变换性质与平行传输}

Wilson线的关键性质是其\textbf{双定域规范变换}行为~\cite{GattringerLang2010}：
\begin{equation}
U_\Gamma(y,x) \;\longrightarrow\; \Omega(y)\, U_\Gamma(y,x)\, \Omega^\dagger(x),
\label{eq:gauge_transform_WL}
\end{equation}
其中$\Omega(x)\in SU(3)$是时空依赖的规范变换矩阵。这一变换性质使Wilson线成为\textbf{颜色空间中的平行传输子}（parallel transporter）——它将费米子场$\psi(x)$从点$x$平行传输到点$y$，使组合$\bar\psi(y)U_\Gamma(y,x)\psi(x)$成为规范不变的。

\textbf{开Wilson线与费米子-反费米子对：}乘积$\bar\psi(y) U_\Gamma(y,x) \psi(x)$在规范变换下的行为可以验证：
\begin{equation}
\bar\psi(y)U_\Gamma(y,x)\psi(x) \to \bar\psi(y)\Omega^\dagger(y) \cdot \Omega(y)U_\Gamma\Omega^\dagger(x) \cdot \Omega(x)\psi(x) = \bar\psi(y)U_\Gamma(y,x)\psi(x),
\end{equation}
因此是一个\textbf{规范不变的夸克-反夸克双定域算符}——这正是所有格点部分子物理的基础。

\subsection{Wilson线作为规范理论的核心范式}

表\ref{tab:WL_roles}总结了Wilson线在规范理论中的四重角色，本文按此框架展开：

\begin{table}[H]
\centering
\caption{Wilson线在格点QCD中的四重角色}
\label{tab:WL_roles}
\small
\begin{tabular}{lp{3cm}p{6cm}}
\toprule
角色 & 涉及的结构 & 物理可观测量 \\
\midrule
\textbf{(A) 禁闭探针} &
矩形Wilson圈、Polyakov圈 &
静态夸克势$V(r)$、弦张力$\sigma$、禁闭-退禁闭相变 \\
\textbf{(B) 格点构建基础} &
规范链接变量$U_\mu(n)$ &
Wilson规范作用、规范不变的费米子作用、haar测度 \\
\textbf{(C) 部分子算符核心} &
直Wilson线、staple型Wilson线 &
准PDF$\tilde q(x,P^z)$、准TMD-PDF$\tilde f(x,b_\perp,P^z)$ \\
\textbf{(D) 非定域探针} &
Landau规范下的长Wilson线 &
规范固定精度检验、重夸克有效理论、红外物理 \\
\bottomrule
\end{tabular}
\end{table}

\subsection{本库中的Wilson线相关文献一览}

\begin{table}[H]
\centering
\caption{本库中Wilson线相关文献的系统分类}
\label{tab:WL_papers}
\small
\begin{tabular}{lp{7cm}}
\toprule
类别 & 关键文献 \\
\midrule
\textbf{教材} &
\texttt{books/Quantum Chromodynamics on the Lattice.pdf} (Ch.2--3,5,7,10)；
\texttt{books/INTRODUCTION TO LATTICE QCD.pdf} (Ch.7,10,14)；

\texttt{books/An Introduction to Quantum Field Theory.pdf} (Part III) \\
\textbf{Wilson线重整化} &
\texttt{Self-Renormalization of Quasi-Light-Front Correlators.pdf}~\cite{Huo2021self}；
\texttt{Improved quasi parton distribution through Wilson line renormalization.pdf}~\cite{Chen2018WLrenorm}；
\texttt{A complete non-perturbative renormalization prescription for quasi-PDFs.pdf} \\
\textbf{TMD的staple Wilson线} &
\texttt{Renormalization of TMD parton distribution on the lattice.pdf}~\cite{Zhang2022renormTMD}；
\texttt{Unpolarized TMD Parton Distributions of the Nucleon from Lattice QCD.pdf}~\cite{He2024unpolTMD}；
\texttt{Quark Transverse Spin-Momentum Correlation...The Boer-Mulders Function.pdf}~\cite{Ma2025BoerMulders} \\
\textbf{准PDF的直Wilson线} &
\texttt{A Hybrid Renormalization Scheme for Quasi Light-Front Correlations.pdf}；
\texttt{Multiplicative renormalizability of quasi-parton operators.pdf}；
\texttt{Resumming Quark's Longitudinal Momentum Logarithms.pdf} \\
\textbf{LaMET因子化} &
\texttt{Factorization Theorem Relating Euclidean and Light-Cone Parton Distributions.pdf}~\cite{Izubuchi2018}；
\texttt{Parton Physics on Euclidean Lattice.pdf}~\cite{Ji2013Euclidean} \\
\textbf{规范固定影响} &
\texttt{Impact of gauge fixing precision on the continuum limit of non-local quark-bilinear lattice operators.pdf}~\cite{Zhang2024gauge_fixing} \\
\bottomrule
\end{tabular}
\end{table}

%==========================================================================
\section{格点上的规范链接：Wilson线的离散化}
%==========================================================================

\subsection{从连续Wilson线到格点规范链接}

在格点上，连续Wilson线被替换为\textbf{规范链接变量}$U_\mu(n)$~\cite{GattringerLang2010}。将连续极限$a\to 0$下的关系展开为：
\begin{equation}
U_\mu(n) = \mathcal{P}\exp\!\left[ig\int_0^a ds\, A_\mu(n a + s\hat\mu)\right] = \exp\!\big[ig a A_\mu(n a) + \mathcal{O}(a^2)\big],
\label{eq:link_variable_def}
\end{equation}
其中$U_\mu(n)$是$SU(3)$群元素，定义在从格点$n$出发、沿$\mu$方向延伸一个格距$a$的\textbf{有向链接}上。$U_{-\mu}(n) \equiv U_\mu(n-\hat\mu)^\dagger$表示沿$-\mu$方向的链接。

\textbf{规范链接的规范变换性质}延续了连续Wilson线的形式：
\begin{equation}
U_\mu(n) \to \Omega(n) U_\mu(n) \Omega^\dagger(n+\hat\mu), \qquad U_{-\mu}(n) \to \Omega(n) U_{-\mu}(n) \Omega^\dagger(n-\hat\mu).
\label{eq:link_gauge_trans}
\end{equation}

\subsection{路径乘积与规范不变量}

任何由规范链接沿路径$\mathcal{P}$的乘积构成的量：
\begin{equation}
P[U] = \prod_{(n,\mu)\in\mathcal{P}} U_\mu(n),
\label{eq:path_product}
\end{equation}
在规范变换下，路径内部的链接矩阵两两相消，仅留下\textbf{两个端点}处的因子：
\begin{equation}
P[U] \to \Omega(n_{\text{start}}) P[U] \Omega^\dagger(n_{\text{end}}).
\label{eq:path_product_gauge}
\end{equation}

由此构造两种基本的\textbf{规范不变量}：

\begin{enumerate}[leftmargin=*]
    \item \textbf{闭合回路（Wilson圈）：}取迹$\operatorname{tr} P[U]$——由于$\operatorname{tr}(AB)=\operatorname{tr}(BA)$，端点矩阵也相消。这是构建规范场作用量和纯规范观测量（如glueball关联函数）的基本范式。
    \item \textbf{费米子封盖的Wilson线：}$\bar\psi(n_{\text{start}})P[U]\psi(n_{\text{end}})$——利用费米子场的规范变换$\psi\to\Omega\psi$、$\bar\psi\to\bar\psi\Omega^\dagger$，端点矩阵与费米子场变换相消。这是所有包含夸克场的非定域算符（如准PDF算符）的基本范式。
\end{enumerate}

\subsection{格点Wilson线的路径求和图像}

Gupta在\textit{Introduction to Lattice QCD}中给出了格点上规范不变量的直观分类~\cite{Gupta1997intro}：
\begin{itemize}
    \item \textbf{开弦（open string）：}路径编序的规范链接链，两端由费米子-反费米子对封盖——在强子关联函数中形成``夸克线'';
    \item \textbf{闭合Wilson圈：}对规范场涨落的最纯粹探针——从$1\times 1$ plaquette到大面积Wilson圈——提取从短距离Coulomb势到长距离线性禁闭势的全部信息。
\end{itemize}

在强耦合展开中，费米子传播子和Wilson圈均可表达为\textbf{规范链接的路径乘积}的级数展开，这是一种``随机曲面''（random surface）的语言——Wilson圈对应于以圈为边界的极小曲面，其面积给出禁闭弦张力的起源。

%==========================================================================
\section{Wilson圈与静态夸克势}
%==========================================================================

\subsection{Wilson圈的构造：矩形时空回路}

\begin{definition}[矩形Wilson圈~\cite{GattringerLang2010}]
考虑一个$r \times t$的矩形时空闭合回路$\mathcal{C}$，由四个部分构成：
\begin{enumerate}[nosep]
    \item 空间Wilson线$S(m, n, t)$——在固定时间片$t$上，从空间点$m$到$n$的规范链接乘积（通常取直线路径）；
    \item 空间Wilson线$S(m, n, 0)$——在时间片$t=0$上的反向传播子（厄米共轭）；
    \item 时间传输子$T(n, t)$——在固定空间点$n$上，沿时间方向的$t$个链接的乘积：$T(n, t) = \prod_{j=0}^{t-1} U_4(n, j)$；
    \item 时间传输子$T(m, t)^\dagger$——在$m$点的反向时间传输。
\end{enumerate}
取迹后，Wilson圈定义为：
\begin{equation}
W_{\mathcal{C}}[U] = \frac{1}{N_c}\operatorname{tr}\!\big[S(m,n,t)\, T(n,t)^\dagger\, S(m,n,0)^\dagger\, T(m,t)\big].
\label{eq:wilson_loop_rectangle}
\end{equation}
\end{definition}

\textbf{平面与非平面Wilson圈：}若$m$和$n$位于同一坐标轴上（空间连线为直路径），则为\textbf{平面}Wilson圈；否则为\textbf{非平面}Wilson圈，后者可用于检验连续极限下旋转不变性的恢复和测量非轴上距离的静态势。

\subsection{Wilson圈的量子力学诠释：从关联函数到势}

利用\textbf{时间规范}$A_4=0$（格点上等价于将时间方向的规范链接全部规范变换为单位矩阵），Wilson圈简化为两个空间Wilson线的关联函数~\cite{GattringerLang2010}：
\begin{equation}
\langle W_{\mathcal{C}} \rangle = \frac{1}{N_c} \langle \operatorname{tr}[S(m,n,t) S(m,n,0)^\dagger] \rangle.
\label{eq:WL_correlator}
\end{equation}

在时间方向的量子力学诠释（式(\ref{eq:quantum_mech})）下，此关联函数在大$t$极限下由\textbf{最低能量本征态}主导：
\begin{equation}
\langle W_{\mathcal{C}} \rangle \;\propto\; e^{-t\,V(r)} \Big[1 + \mathcal{O}(e^{-t\Delta E})\Big],
\label{eq:WL_potential}
\end{equation}
其中$r = a|m-n|$是夸克-反夸克的空间距离，$V(r)$是\textbf{静态夸克-反夸克势}——包含Coulomb部分和线性禁闭部分：
\begin{equation}
\boxed{V(r) = A + \frac{B}{r} + \sigma r},
\label{eq:static_potential}
\end{equation}
$A$是自能常数（不可观测），$B$是Coulomb系数，$\sigma$是\textbf{弦张力}（$\sigma \approx 900$~MeV/fm $\approx 4.5$~MeV/fm$^2\,$?）。

\textbf{Wilson线的变换性质验证：}在Chap.~5（费米子）中，Gattringer与Lang~\cite{GattringerLang2010}通过hopping展开证明：在无穷大夸克质量（静态）极限下，费米子传播子退化为空间Wilson线$S(m,n)$——这确证了Wilson线确实描述\textbf{静态夸克-反夸克对}的状态。

\subsection{强耦合展开与面积律：禁闭的格点证明}

强耦合（小$\beta = 6/g^2$）展开是格点QCD中证明线性禁闭的经典方法~\cite{GattringerLang2010,Wilson1974}。

\textbf{核心计算（式(3.68)）：}Wilson圈期望值在强耦合下的领头阶贡献为：
\begin{equation}
\langle W_{\mathcal{C}} \rangle = 3 \exp\!\left[ n_r n_t \ln\frac{\beta}{18} \right] (1+\mathcal{O}(\beta)) = 3 \exp\!\left[ -\sigma a^2 \cdot n_r n_t \right],
\label{eq:strong_coupling_WL}
\end{equation}
其中$n_r n_t$是矩形Wilson圈包围的\textbf{最小面积}（以plaquette计）。由此导出弦张力的强耦合表达式：
\begin{equation}
\sigma = -\frac{1}{a^2}\ln\frac{\beta}{18} + \mathcal{O}(\beta).
\label{eq:string_tension_SC}
\end{equation}

\textbf{面积律的物理图像（图3.4）：}Wilson圈的迹$\operatorname{tr}[\prod U_l]$中的每个规范链接需要被玻尔兹曼因子展开中的\underline{反向plaquette}来``配对''——形成以Wilson圈为边界、plaquette为瓦片的\textbf{极小曲面}。曲面面积$n_r n_t$越大，需要的plaquette数目越多，贡献的$\beta$幂次越高，因此期望值按面积指数衰减——这便是\textbf{面积律}（area law）的格点起源。

\subsection{Sommer参数与格距标定}

从Wilson圈提取出的静态势$V(r)$不仅可以证明禁闭，还可以用于\textbf{标定格距}~\cite{GattringerLang2010}。Sommer参数$r_0$通过力的条件来定义：
\begin{equation}
F(r_0) r_0^2 = 1.65, \qquad r_0 \simeq 0.5\text{~fm},
\label{eq:sommer_parameter}
\end{equation}
其中$F(r) = dV/dr$。从格点数据中确定$r_0/a$（Sommer尺度在格点单位下的值），即可标定格距$a$。

\textbf{本库中的相关应用：}在LPC合作组的TMD-PDF计算~\cite{He2024unpolTMD}中，MILC系综的格距$a=0.12$~fm正是通过类似的标定方法确定的。CLS系综的$a=0.098$~fm~\cite{Chu2023soft}同样依赖于高精度静态势测量。

\subsection{Wilson线的静态极限与重夸克有效理论的联系}

Wilson线作为静态夸克传播子的物理解释直接链接到\textbf{重夸克有效理论（HQET）}。在$m_Q\to\infty$极限下，重夸克的传播子退化为Wilson线：
\begin{equation}
\langle 0| Q(x) \bar Q(0) |0 \rangle \;\xrightarrow{m_Q\to\infty}\; U(x,0),
\end{equation}
其中$U(x,0)$是沿时间（或空间）方向连接两点的Wilson线。在HQET的格点实现中，静态重夸克作用量直接由时间方向的Wilson线构建~\cite{GattringerLang2010}。

这一联系对于理解部分子物理中的准PDF算符至关重要——准PDF算符中的空间方向Wilson线可以视为``重夸克化''后的规范平行传输子，其线性发散的结构与HQET中的重夸克自能修正完全一致。

%==========================================================================
\section{Polyakov圈与有限温度物理}
%==========================================================================

\subsection{Polyakov圈：周期边界条件驱动的热力学Wilson线}

在有限温度格点QCD中（时间方向取周期边界条件，长度$\beta = 1/T = aN_T$），自然产生一类特殊的闭合Wilson线——Polyakov圈~\cite{GattringerLang2010,Gupta1997intro}：

\begin{definition}[Polyakov圈（热力学Wilson线）]
\begin{equation}
P(\mathbf{m}) \equiv \operatorname{tr}\!\left[ \prod_{j=0}^{N_T-1} U_4(\mathbf{m}, j) \right],
\label{eq:polyakov_loop}
\end{equation}
即\textbf{纯时间方向}、贯穿整个时间范围$[0, N_T-1]$的闭合Wilson线——由于周期边界条件，时间方向首尾自动相连形成闭合回路。
\end{definition}

\textbf{与Wilson圈的关系：}令矩形Wilson圈的时间跨度$n_t = N_T$（最大可能值），则在适当的规范下，空间部分退耦，Wilson圈简化为两个Polyakov圈$P(\mathbf{m})$和$P(\mathbf{n})^\dagger$的乘积。Polyakov圈关联函数同样给出静态势：
\begin{equation}
\langle P(\mathbf{m}) P(\mathbf{n})^\dagger \rangle \propto e^{-N_T a V(r)}, \qquad r = a|\mathbf{m}-\mathbf{n}|.
\label{eq:polyakov_correlator}
\end{equation}

\subsection{中心对称性与禁闭-退禁闭相变}

Polyakov圈是有限温度下\textbf{禁闭-退禁闭相变的有序参量}~\cite{GattringerLang2010}。纯规范理论中，存在一个$Z_3$中心对称性：将某个时间片上的所有时间方向链接乘以$z\in Z_3$（$1, e^{2\pi i/3}, e^{-2\pi i/3}$），规范作用保持不变。Polyakov圈在此变换下获得因子$z$：
\begin{equation}
P \to zP.
\label{eq:center_transform}
\end{equation}

在\textbf{禁闭相}（$T < T_c \approx 270$~MeV），中心对称性保持，$Z_3$对称性迫使$\langle P \rangle = 0$——单个静态色电荷的自由能为无穷大。
在\textbf{退禁闭相}（$T > T_c$），中心对称性自发破缺，$\langle P \rangle \neq 0$——色电荷被Debye屏蔽，自由能有限。

\textbf{动力学费米子的效应：}动力学费米子显式破缺中心对称性——费米子行列式中包含Polyakov圈的贡献（$\kappa^{N_T}(P+P^*)$），类似于自旋模型中的外磁场。因此，在full QCD中，Polyakov圈不是严格的有序参量，禁闭-退禁闭转变为平滑的\textbf{crossover}~\cite{GattringerLang2010}。

%==========================================================================
\section{部分子物理中的Wilson线：从准PDF到准TMD-PDF}
%==========================================================================

\subsection{直Wilson线与准PDF算符}

在LaMET框架~\cite{Ji2013Euclidean,Ji2014LaMET}中，共线PDF的格点计算始于\textbf{准PDF算符}——一个由空间方向直Wilson线连接的非定域夸克双线型：

\begin{definition}[准PDF算符中的直Wilson线]
对于沿$z$方向boost的强子，等时准PDF算符为：
\begin{equation}
O_\Gamma(z) \equiv \bar\psi(z\hat n_z) \Gamma U_z(z,0) \psi(0),
\label{eq:quasi_PDF_operator}
\end{equation}
其中\textbf{直Wilson线}（straight Wilson line）$U_z(z,0)$沿$z$轴连接夸克场：
\begin{equation}
U_z(z,0) \equiv \mathcal{P}\exp\!\left[ig\int_0^z ds\, A_z(s\hat n_z)\right] \quad\longrightarrow\quad \prod_{j=0}^{z/a-1} U_z(j\,\hat n_z) \;\;(\text{格点上}).
\label{eq:straight_WL}
\end{equation}
\end{definition}

\textbf{物理意义：}在光锥坐标系中，光锥PDF定义为沿光锥方向$n_-$的夸克-反夸克关联函数。LaMET的核心洞察是：在大动量$P^z$下，等时空间方向的关联函数可以通过微扰匹配与光锥关联函数联系起来。空间方向的直Wilson线$U_z$正是光锥方向上规范链接的等时对应物。

\textbf{准PDF的完整定义~\cite{Ji2013Euclidean}：}
\begin{equation}
\tilde q(x, P^z) \equiv \int \frac{dz}{2\pi} e^{-i x P^z z} \langle P^z | \bar\psi(z\hat n_z) \gamma^t U_z(z,0) \psi(0) | P^z \rangle.
\label{eq:quasi_PDF_full}
\end{equation}

\subsection{Staple型Wilson线与准TMD-PDF算符}

TMD-PDF的描述需要同时捕获\textbf{纵向}和\textbf{横向}的夸克动量分布，因此其准TMD-PDF算符需要更复杂的Wilson线结构~\cite{Zhang2022renormTMD,He2024unpolTMD,Ma2025BoerMulders}。

\begin{definition}[准TMD-PDF的Staple型Wilson线~\cite{Zhang2022renormTMD}]
连接夸克场$\bar\psi(\vec{b}_\perp, 0)$和$\psi(z\hat n_z)$的staple型规范链接由三段Wilson线组成：
\begin{equation}
\begin{aligned}
W_{\sqsubset}(b_\perp, L, z) &\equiv U_z^\dagger\big((z+L)\hat n_z + \vec{b}_\perp,\; \vec{b}_\perp\big) \\
&\quad \times U_\perp\big((z+L)\hat n_z + \vec{b}_\perp,\; (z+L)\hat n_z\big) \\
&\quad \times U_z\big((z+L)\hat n_z,\; z\hat n_z\big),
\end{aligned}
\label{eq:staple_WL}
\end{equation}
其中$b_\perp = |\vec{b}_\perp|$是横向分离，$z$是纵向分离，$L$是staple的``臂长''——向$z$方向延伸到``无穷远''（格点上为有限值，随后外推$L\to\infty$）。
\end{definition}

\textbf{Staple的三段Wilson线各有物理功能：}
\begin{enumerate}[nosep]
    \item $U_z((z+L)\hat n_z, z\hat n_z)$——从$z$沿$z$正向延伸到$z+L$，为纵向传输；
    \item $U_\perp$——在``无穷远''处（$z+L$）进行横向跳跃$\vec{b}_\perp$，使Wilson线跨越到横向偏移的端点；
    \item $U_z^\dagger((z+L)\hat n_z+\vec{b}_\perp, \vec{b}_\perp)$——从$\vec{b}_\perp$反向沿$z$延伸到$z+L+\vec{b}_\perp$，完成到横向位移点的传输。
\end{enumerate}
这一结构与TMD因子化中\textbf{两个相反光锥方向上的Wilson线}的连续极限结构完全对应。

\textbf{准TMD-PDF的裸矩阵元~\cite{Zhang2022renormTMD}：}
\begin{equation}
\tilde{h}_{\chi,\Gamma}(b_\perp, z, L, P^z; 1/a) = \langle \chi(P^z) | \bar\psi(\vec{0}_\perp, 0) \Gamma W_{\sqsubset}(b_\perp, L, z) \psi(\vec{b}_\perp, z) | \chi(P^z) \rangle.
\label{eq:bare_TMD_matrix}
\end{equation}

对于不同的TMD-PDF，$\Gamma$矩阵的选取不同：非极化$f_1$取$\Gamma = \gamma^t$，Boer-Mulders函数$h_1^\perp$取$\Gamma = i\sigma^{yt}\gamma_5$~\cite{Ma2025BoerMulders}。

\subsection{Wilson线在因子化定理中的核心地位}

Izubuchi等人的因子化定理~\cite{Izubuchi2018}确立了准TMD-PDF与光锥TMD-PDF之间的精确对应关系。在此定理中，软函数$S(b_\perp, \mu, Y+Y')$——它描述了来自两个相反光锥方向Wilson线的软胶子辐射——起着中心作用。软函数本身直接由真空中的staple型Wilson圈定义，是Wilson线物理在TMD因子化中的最深远体现。

%==========================================================================
\section{Wilson线的发散结构与重整化}
%==========================================================================

\subsection{线性发散的起源与结构}

Wilson线的\textbf{自能修正}产生格点QCD中最棘手的一类紫外发散——\textbf{线性发散}（power divergence）~\cite{Huo2021self,Chen2018WLrenorm}。考虑一条长度为$L$的Wilson线（格点上为$L/a$个链接），其裸期望值在格点正规化下具有以下结构：

\begin{equation}
\langle U(L,0) \rangle_{\text{bare}} \sim e^{-\delta\bar{m} L/a} \times (\text{对数发散因子}) \times (\text{物理贡献}),
\label{eq:WL_bare_structure}
\end{equation}
其中$\delta\bar{m} \propto g^2$是\textbf{线性发散系数}（``残余质量''）——具有质量量纲，在连续极限$a\to 0$下发散。

\textbf{线性发散的物理本质：}在微扰论中，线性发散来自Wilson线自能图（蝌蚪图）中胶子传播子在UV区域的行为。在辅助$z$场形式中~\cite{Chen2018WLrenorm}，Wilson线等价于一个具有``重夸克''传播子特征的非动力学场$Z(z)$的二点函数。线性发散的结构与重夸克有效理论（HQET）中的\textbf{残余质量项}$\delta m \,\bar h_v h_v$完全相同——两者都需要一个\textbf{质量抵消项}来吸收。

\textbf{线性发散的$z$依赖性：}在准PDF算符$\bar\psi(z)\Gamma U(z,0)\psi(0)$的矩阵元中，线性发散出现在指数因子里：$e^{-\delta\bar{m} |z|/a}$。这意味着矩阵元随$|z|$增大而指数衰减——即使物理的共线PDF并没有这样的衰减。在连续极限$a\to 0$下，固定物理$z$的衰减速率发散，使得裸矩阵元\textbf{不直接具有连续极限}。

\subsection{策略I：质量抵消项与辅助$z$场形式}

Chen等人~\cite{Chen2018WLrenorm}证明：线性发散可以通过\textbf{质量抵消项}$\delta m$来系统地消除——这是对开放Wilson线的所有阶都成立的结论。

\textbf{辅助$z$场构造：}将Wilson线$U(z,0)$视为连接$Z(0)$和$\bar Z(z)$两个辅助场的传播子：
\begin{equation}
U(z,0) \sim \langle 0| Z(z) \bar Z(0) |0 \rangle,
\end{equation}
其中$Z(z)$场在路径积分中不具有动力学动能项（在HQET极限下）。在此形式下，线性发散对应于$Z$场的质量重整化，可以通过标准的\textbf{在壳质量抵消项}$\delta m$来消除。

\textbf{到所有阶的证明：}任何包含Wilson线的Feynman图，其线性发散仅来源于Wilson线自能子图。质量抵消项$\delta m$在这些子图中与线性发散精确相消——这一抵消在所有微扰阶均成立，因为线性发散的系数$\delta\bar{m}$具有普适性（不依赖于具体的夸克双线型或外部态）。

\subsection{策略II：Wilson圈减除——TMD-PDF的标准方案}

LPC合作组~\cite{Zhang2022renormTMD}开发了针对准TMD-PDF的\textbf{Wilson圈减除}重整化方案，可同时消除三类发散。

\begin{definition}[Wilson圈减除的准TMD-PDF~\cite{Zhang2022renormTMD}]
\begin{equation}
\boxed{h_{\chi,\Gamma}(b_\perp, z, P^z; 1/a) \equiv \lim_{L\to\infty} \frac{\tilde{h}_{\chi,\Gamma}(b_\perp, z, L, P^z; 1/a)}{\sqrt{Z_E(2L+z, b_\perp; 1/a)}}},
\label{eq:WL_subtraction}
\end{equation}
其中\textbf{矩形Wilson圈}$Z_E$在真空中的定义为：
\begin{equation}
\begin{aligned}
Z_E(r, b_\perp; 1/a) &\equiv \frac{1}{N_c} \operatorname{Tr} \langle 0| U_\perp^\dagger(-\vec{L}+\vec{b}_\perp; -b_\perp) U_z^\dagger(\vec{L}+\vec{z}+\vec{b}_\perp; -2L-z) \\
&\qquad \times U_\perp(\vec{L}+\vec{z}; b_\perp) U_z(-\vec{L}; 2L+z) |0\rangle.
\label{eq:ZE_wilson_loop}
\end{aligned}
\end{equation}
\end{definition}

\textbf{减除的物理机制——三种发散的消除：}
\begin{enumerate}[leftmargin=*]
    \item \textbf{线性发散：}分母$\sqrt{Z_E}$中的Wilson线总长同样为$2L+z$，因此其指数因子为$e^{-\delta\bar{m}(2L+z)/a}$。开方后恰好消除分子的线性发散$e^{-\delta\bar{m}(2L+z)/a}$。
    \item \textbf{Pinch-pole奇点：}Staple中两条平行$z$方向Wilson线之间的胶子交换产生$e^{-V(b_\perp)L}$因子（物理上等价于重夸克势）。矩形Wilson圈$Z_E$中的两对平行边产生相同的奇点——开方抵消。
    \item \textbf{Cusp发散：}Staple的三个转角处的尖角发散和矩形Wilson圈的四个转角处的尖角发散在减除中消除。
\end{enumerate}

\textbf{五格点间距的系统验证~\cite{Zhang2022renormTMD}：}Zhang等人在$a = 0.032, 0.043, 0.060, 0.086, 0.121$~fm上系统验证：
\begin{itemize}
    \item 减除后的矩阵元在$L \gtrsim 0.36$~fm后达到\textbf{平台}——$L$独立性验证了pinch-pole消除；
    \item 在$a\to 0$连续极限下稳定收敛——验证了线性发散的消除；
    \item Clover和Overlap费米子结果一致——确认了重整化的\textbf{作用量普适性}。
\end{itemize}

\textbf{减除后残余的对数发散}通过\textbf{短距离比率方案}（SDR）处理——利用零动量和微扰短距离处的矩阵元形成比率，将格点正规化匹配到$\overline{\text{MS}}$方案~\cite{Zhang2022renormTMD,He2024unpolTMD}。

\subsection{策略III：自重整化方法}

Huo等人~\cite{Huo2021self}提出了\textbf{自重整化}（self-renormalization）方法——一个不依赖于微扰计算的纯非微扰方案。

\begin{definition}[自重整化策略]
核心思想是利用同一矩阵元在不同$z$处的比值消除线性发散和重整子（renormalon）效应：
\begin{equation}
M^{\text{ren}}(z, P^z) = \frac{M^{\text{bare}}(z, P^z)}{\big[M^{\text{bare}}(z_0, P^z=0)\big]^{z/z_0}},
\label{eq:self_renorm}
\end{equation}
其中$z_0$是一个固定的短距离参考点。由于线性发散因子$e^{-\delta\bar{m}|z|/a}$在比值中的指数结构被约化为1，所有与$z$成正比的发散被自动消除。
\end{definition}

\textbf{自重整化的优势：}
\begin{itemize}
    \item 纯非微扰——不需要Wilons圈测量，不需要微扰匹配；
    \item 适用于直Wilson线（准PDF）和staple Wilson线（准TMD-PDF）；
    \item 可同时消除线性发散和重整子模糊（renormalon ambiguity）；
    \item 对HYP涂抹前后的数据均有效。
\end{itemize}

\textbf{方法验证~\cite{Huo2021self}：}在Clover和Overlap费米子的多个系综上，自重整化后的矩阵元在不同格距上一致——确认了消除线性发散的普适性。

\subsection{RI/MOM方案的失败}

Zhang等人~\cite{Zhang2022renormTMD}和Huo等人~\cite{Huo2021self}一致发现：传统的RI/MOM方案\textbf{无法完全消除}准TMD-PDF和准PDF算符中的线性发散。

\textbf{失败的表现：}RI/MOM重整化后的矩阵元在更大$z$（或更小$a$）时展示\textbf{更强的格距依赖性}——意味着$a\to 0$极限不收敛。

\textbf{根本原因：}在破坏规范对称性的格点截断方案下，胶子-规范链接（Wilson线）顶点处的线性发散是非微扰本质的，而RI/MOM方案（基于离壳格林函数在微扰标度$\mu$处的减除）无法捕获这种非微扰线性发散。在Overlap费米子（保持手征对称性）上的交叉验证证实了这是一种\textbf{普遍性缺陷}，而非特定费米子作用量的artifacts。

\textbf{核心教训：}对于包含Wilson线的非定域算符，\textbf{基于减除的方案}（Wilson圈、自重整化、或混合方案中的ratio部分）是唯一可行的重整化路径。

%==========================================================================
\section{Wilson线的改进：HYP涂抹与规范固定}
%==========================================================================

\subsection{HYP涂抹：改善Wilson线的信噪比}

HYP（Hypercubic）涂抹~\cite{HYPsmearing}是一种广泛应用于格点QCD的规范链接平滑技术，其动机是改善包含Wilson线的非定域算符的\textbf{信噪比}。

\textbf{HYP涂抹的构造：}在hypercubic的一个子集内，目标链接$U_\mu(n)$被替换为其周围链接的加权平均（投影到$SU(3)$后的结果）。与传统的APE或stout涂抹不同，HYP涂抹在三个正交的hypercubic平面内分别进行平滑，以保持尽可能高的局域性。

\textbf{HYP涂抹对Wilson线的影响~\cite{Huo2021self}：}
\begin{itemize}
    \item 涂抹显著降低了Wilson线自能的紫外涨落，从而\textbf{压制了线性发散}的绝对大小（$\delta\bar{m}$在涂抹后显著减小）——但需注意涂抹后的线性发散虽减小，却未被完全消除；
    \item 涂抹也改变了线性发散系数的非微扰值，使其对格点作用量和费米子方案的依赖性减弱；
    \item 涂抹后的Wilson线在自重整化框架下同样适用，但\underline{直接比较}涂抹前后的重整化矩阵元需要谨慎——涂抹改变了算符的UV结构。
\end{itemize}

\textbf{本库中的应用：}LPC合作组的TMD-PDF计算~\cite{He2024unpolTMD}中，规范链接均使用了HYP涂抹，显著改善了核子三点函数的信噪比。

\subsection{规范固定精度对长Wilson线的影响}

Zhang等人~\cite{Zhang2024gauge_fixing}系统研究了Landau规范固定精度对非定域夸克双线型算符（即包含Wilson线的算符）连续极限行为的影响。

\textbf{核心发现：}包含Wilson线的算符对规范固定的精度\textbf{极为敏感}，远超之前的预期：
\begin{itemize}
    \item 在$a = 0.03$~fm的细格距上，即使使用通常认为``足够''的规范固定精度$\delta_F = 10^{-8}$，1.2~fm长的Wilson线的迹仍可能产生约\textbf{12\%}的偏差；
    \item 影响随Wilson线长度增加呈\textbf{二次方增长}：$c(W(z)) \propto z^2$，意味着长距离下的规范固定不精确性是一个系统性的限制因素；
    \item 细格距（小$a$）上的影响比粗格距更严重——这是因为细格距上格点数量更大，规范固定需要收敛的自由度更多。
\end{itemize}

\textbf{对实践的启示：}在进行包含长Wilson线的计算（如大$x$区域的准PDF或大$b_\perp$的准TMD-PDF）时，需要将规范固定精度推至远高于常规水平（$\delta_F \lesssim 10^{-12}$可能对细格距是必要的）。这为本库中的准PDF和准TMD-PDF计算提供了重要的\textbf{系统误差控制指南}。

%==========================================================================
\section{总结与展望}
%==========================================================================

基于对三本核心教材和约35篇研究论文的综合分析，本报告从四个层面系统阐述格点QCD中Wilson线的理论、计算与物理：

\begin{enumerate}[leftmargin=*]

    \item \textbf{Wilson线是格点QCD的``基本砖块''：}从连续路径编序指数$U_\Gamma = \mathcal{P}\exp(ig\int A_\mu dx^\mu)$到格点规范链接$U_\mu(n)$，Wilson线以平行传输子的身份构建了理论的规范结构。闭合Wilson圈和费米子封盖的开Wilson线是仅有的两类规范不变量——前者定义纯规范作用量和glueball谱，后者是所有包含夸克的非定域算符的范式。

    \item \textbf{Wilson圈给出了禁闭的格点证明：}面积律$\langle W_{\mathcal{C}}\rangle \propto \exp(-\sigma\cdot\text{Area})$通过强耦合展开从第一性原理导出，确证了线性禁闭势$V(r) \sim \sigma r$的存在。Sommer参数$r_0$使得格距可以从Wilson圈的数值数据中标定。Polyakov圈作为热力学Wilson线，在纯规范理论中充当禁闭-退禁闭相变的严格有序参量。

    \item \textbf{现代部分子物理以Wilson线为核心算符要素：}LaMET框架将光锥关联函数替换为等时空间关联函数，其中直Wilson线（准PDF）和staple型Wilson线（准TMD-PDF）是算符不可分割的组成部分。从Izubuchi等人的因子化定理~\cite{Izubuchi2018}到He等人的核子TMD-PDF首次计算~\cite{He2024unpolTMD}，Wilson线的理论处理能力直接决定了部分子物理格点计算的精度边界。

    \item \textbf{线性发散的重整化是Wilson线物理的核心技术挑战：}三种方案——质量抵消项（Chen等人~\cite{Chen2018WLrenorm}）、Wilson圈减除（Zhang等人~\cite{Zhang2022renormTMD}）和自重整化（Huo等人~\cite{Huo2021self}）——均已被证明可在不同场景下有效消除线性发散。RI/MOM方案的失败~\cite{Zhang2022renormTMD,Huo2021self}确立了减除方案的唯一性。HYP涂抹可压制但非消除线性发散。规范固定精度对长Wilson线的显著影响~\cite{Zhang2024gauge_fixing}揭示了精细格距计算中新的系统误差来源。

    \item \textbf{Wilson线将格点QCD的三个尺度——禁闭（$\sim$ fm以上）、强子结构（$\sim 0.1-1$~fm）和微扰匹配（$\lesssim 0.1$~fm）——编织成一个统一的计算框架，}串联起从弦张力提取、格距标定、PDF/TMD-PDF计算到连续极限外推的全链条。

\end{enumerate}

\textbf{开放问题与展望：}
\begin{itemize}
    \item Staple型Wilson线的非微扰重整化是否可推广到胶子TMD算符（需伴随表示Wilson线）？
    \item Wilson线减除方案在更精细的格距（$a \lesssim 0.03$~fm）和更长的Wilson线（$z \gtrsim 1.5$~fm）下是否仍保持有效性？
    \item 自重整化方法和Wilson圈减除方法在共同的连续极限下是否给出相同的结果？二者之间的系统偏差需定量评估。
\end{itemize}

%==========================================================================
\section*{本库中Wilson线相关文献的全目录}
%==========================================================================

\small
\begin{enumerate}[leftmargin=*,nosep]
    \item \texttt{books/Quantum Chromodynamics on the Lattice.pdf} [Gattringer \& Lang, 2010] — Ch.2 (规范链接), Ch.3.3--3.5 (Wilson圈、静态势、Sommer参数), Ch.5 (hopping展开与Wilson线的静态夸克解释), Ch.7 (手征对称性与Ginsparg--Wilson费米子), Ch.10 (Staggered/Domain-Wall/Twisted Mass费米子)
    \item \texttt{books/INTRODUCTION TO LATTICE QCD.pdf} [R. Gupta, 1997] — Ch.7 (费米子作用), Ch.10 (Wilson费米子), Ch.14 (Wilson圈、弦张力)
    \item \texttt{books/An Introduction to Quantum Field Theory.pdf} [Peskin \& Schroeder, 1995] — Part III (连续QCD背景)
    \item \texttt{docs/Self-Renormalization of Quasi-Light-Front Correlators on the Lattice.pdf}~\cite{Huo2021self} — 自重整化方法、线性发散分析在Clover/Overlap费米子上的系统验证
    \item \texttt{docs/Improved quasi parton distribution through Wilson line renormalization.pdf}~\cite{Chen2018WLrenorm} — 辅助$z$场形式、质量抵消项的系统证明
    \item \texttt{docs/Renormalization of transverse-momentum-dependent parton distribution on the lattice.pdf}~\cite{Zhang2022renormTMD} — 五种格距间距上Wilson圈减除的系统验证、RI/MOM失败证明
    \item \texttt{docs/Impact of gauge fixing precision on the continuum limit of non-local quark-bilinear lattice operators.pdf}~\cite{Zhang2024gauge_fixing} — 规范固定精度对长Wilson线的定量影响
    \item \texttt{docs/Unpolarized Transverse-Momentum-Dependent Parton Distributions of the Nucleon from Lattice QCD.pdf}~\cite{He2024unpolTMD} — 核子非极化TMD-PDF的首次计算（含staple Wilson线+HYP涂抹）
    \item \texttt{docs/Quark Transverse Spin-Momentum Correlation of the Nucleon from Lattice QCD The Boer-Mulders Function.pdf}~\cite{Ma2025BoerMulders} — Boer-Mulders函数的首次核子计算（含staple Wilson线）
    \item \texttt{docs/Parton Physics on Euclidean Lattice.pdf}~\cite{Ji2013Euclidean} — LaMET的奠基性论文，直Wilson线和staple Wilson线的原始引入
    \item \texttt{docs/Factorization Theorem Relating Euclidean and Light-Cone Parton Distributions.pdf}~\cite{Izubuchi2018} — 准TMD-PDF与光锥TMD-PDF之间因子化定理的严格证明
    \item \texttt{docs/Parton Physics from Large-Momentum Effective Field Theory.pdf}~\cite{Ji2014LaMET} — LaMET框架的系统阐述
    \item \texttt{docs/A complete non-perturbative renormalization prescription for quasi-PDFs.pdf} — 准PDF的完整非微扰重整化方案
    \item \texttt{docs/A Hybrid Renormalization Scheme for Quasi Light-Front Correlations in Large-Momentum Effective Theory.pdf} — 混合重整化方案
    \item \texttt{docs/Multiplicative renormalizability of quasi-parton operators.pdf} — 准部分子算符的乘法重整化性
    \item \texttt{docs/Transverse-Momentum-Dependent Wave Functions of Pion from Lattice QCD.pdf}~\cite{Chu2023TMDWF} — $\pi$介子TMD波函数计算（含staple Wilson线）
    \item \texttt{docs/Lattice Calculation of the Intrinsic Soft Function and the Collins-Soper Kernel.pdf}~\cite{Chu2023soft} — 软函数计算（含真空Wilson圈）
    \item \texttt{docs/Nonperturbative Determination of Collins-Soper Kernel from Quasi Transverse-Momentum Dependent Wave Functions.pdf} — CS核的准TMD波函数提取
    \item \texttt{docs/Lattice-QCD Calculations of TMD Soft Function Through Large-Momentum Effective Theory.pdf}~\cite{LPC2020soft} — 内禀软函数和CS核的首次格点计算
    \item \texttt{docs/Gauge-Invariant Smearing and Matrix Correlators using Wilson Fermions at beta=6.2.pdf} — Wilson费米子上的规范不变量涂抹（Wilson圈测量技术）
\end{enumerate}

%==========================================================================
\begin{thebibliography}{99}

\bibitem{GattringerLang2010}
C.~Gattringer and C.~B.~Lang,
\textit{Quantum Chromodynamics on the Lattice: An Introductory Presentation},
Lect.\ Notes Phys.\ \textbf{788}, Springer (2010).
\\\textbf{[来源:} \texttt{books/Quantum Chromodynamics on the Lattice.pdf}\textbf{]}
\\Ch.2--3（规范链接、Wilson圈、静态势）、Ch.5（hopping展开与Wilson线的静态夸克解释）、Ch.7（Ginsparg--Wilson关系/Overlap）、Ch.9（Symanzik改进/Clover）、Ch.10（Staggered/Domain-Wall/Twisted Mass）。

\bibitem{Gupta1997intro}
R.~Gupta,
``Introduction to Lattice QCD,''
arXiv:hep-lat/9807028 (1997).
\\\textbf{[来源:} \texttt{books/INTRODUCTION TO LATTICE QCD.pdf}\textbf{]}
\\Ch.7（费米子作用中的规范链接）、Ch.10（Wilson费米子）、Ch.14（Wilson圈、弦张力、Polyakov圈）。

\bibitem{PeskinSchroeder1995}
M.~E.~Peskin and D.~V.~Schroeder,
\textit{An Introduction to Quantum Field Theory},
Westview Press (1995).
\\\textbf{[来源:} \texttt{books/An Introduction to Quantum Field Theory.pdf}\textbf{]}

\bibitem{Wilson1974}
K.~G.~Wilson,
``Confinement of quarks,''
Phys.\ Rev.\ D \textbf{10}, 2445 (1974).
——Wilson圈作为禁闭探针的原始提出，包括面积律的强耦合推导和Wilson费米子作用。

\bibitem{Huo2021self}
Y.-K.~Huo \textit{et al.} (LPC),
``Self-renormalization of quasi-light-front correlators on the lattice,''
Nucl.\ Phys.\ B \textbf{969}, 115443 (2021), arXiv:2103.02965.
\\\textbf{[来源:} \texttt{docs/Self-Renormalization of Quasi-Light-Front Correlators on the Lattice.pdf}\textbf{]}
——自重整化方法：通过同一矩阵元在不同$z$处的比值消除线性发散和重整子模糊。
Clover和Overlap费米子上对直Wilson线和HYP涂抹数据的系统验证（本报告第6.3节核心参照）。

\bibitem{Chen2018WLrenorm}
J.-W.~Chen, T.~Ishikawa, L.~Jin, H.-W.~Lin, Y.-B.~Yang, J.-H.~Zhang, and Y.~Zhao,
``Improved quasi parton distribution through Wilson line renormalization,''
Phys.\ Rev.\ D \textbf{97}, 014505 (2018), arXiv:1706.01295.
\\\textbf{[来源:} \texttt{docs/Improved quasi parton distribution through Wilson line renormalization.pdf}\textbf{]}
——辅助$z$场形式中线性发散作为质量重整化的证明，质量抵消项$\delta m$到所有阶的有效性。

\bibitem{Zhang2022renormTMD}
K.~Zhang \textit{et al.} (LPC),
``Renormalization of transverse-momentum-dependent parton distribution on the lattice,''
Phys.\ Rev.\ D \textbf{106}, 094510 (2022), arXiv:2205.13402.
\\\textbf{[来源:} \texttt{docs/Renormalization of transverse-momentum-dependent parton distribution on the lattice.pdf}\textbf{]}
——五种格点间距上Wilson圈减除（式\ref{eq:WL_subtraction}）的系统验证、RI/MOM失败证明。
Clover和Overlap费米子的交叉验证（本报告第6.2和6.4节核心参照）。

\bibitem{Zhang2024gauge_fixing}
K.~Zhang \textit{et al.} (LPC),
``Impact of gauge fixing precision on the continuum limit of non-local quark-bilinear lattice operators,''
Phys.\ Rev.\ D \textbf{110}, 074505 (2024), arXiv:2405.14097.
\\\textbf{[来源:} \texttt{docs/Impact of gauge fixing precision on the continuum limit of non-local quark-bilinear lattice operators.pdf}\textbf{]}
——规范固定精度对长Wilson线的定量影响：$\delta_F=10^{-8}$，$a=0.03$~fm，$z=1.2$~fm $\to$ 12\%偏差（本报告第7.2节核心参照）。

\bibitem{He2024unpolTMD}
J.-C.~He \textit{et al.} (LPC),
``Unpolarized transverse-momentum-dependent parton distributions of the nucleon from lattice QCD,''
Phys.\ Rev.\ D \textbf{109}, 114513 (2024), arXiv:2211.02340.
\\\textbf{[来源:} \texttt{docs/Unpolarized Transverse-Momentum-Dependent Parton Distributions of the Nucleon from Lattice QCD.pdf}\textbf{]}
——核子非极化TMD-PDF的首次格点计算（含staple Wilson线+HYP涂抹+Wilson圈减除）。

\bibitem{Ma2025BoerMulders}
L.~Ma \textit{et al.} (LPC),
``Quark transverse spin-momentum correlation of the nucleon from lattice QCD: The Boer-Mulders function,''
(2025), arXiv:2502.11807.
\\\textbf{[来源:} \texttt{docs/Quark Transverse Spin-Momentum Correlation of the Nucleon from Lattice QCD The Boer-Mulders Function.pdf}\textbf{]}
——Boer-Mulders函数的首次格点计算（含staple Wilson线，算符$i\sigma^{yt}\gamma_5 W_{\sqsubset}$）。

\bibitem{Ji2013Euclidean}
X.~Ji,
``Parton physics on a Euclidean lattice,''
Phys.\ Rev.\ Lett.\ \textbf{110}, 262002 (2013), arXiv:1305.1539.
\\\textbf{[来源:} \texttt{docs/Parton Physics on Euclidean Lattice.pdf}\textbf{]}
——LaMET的奠基性论文：首次提出利用空间方向直Wilson线和staple型Wilson线构建准PDF和准TMD-PDF算符。

\bibitem{Ji2014LaMET}
X.~Ji,
``Parton physics from large-momentum effective field theory,''
Sci.\ China Phys.\ Mech.\ Astron.\ \textbf{57}, 1407 (2014), arXiv:1404.6680.
\\\textbf{[来源:} \texttt{docs/Parton Physics from Large-Momentum Effective Field Theory.pdf}\textbf{]}
——LaMET框架的系统阐述。

\bibitem{Izubuchi2018}
T.~Izubuchi, X.~Ji, L.~Jin, I.~W.~Stewart, and Y.~Zhao,
``Factorization theorem relating Euclidean and light-cone parton distributions,''
Phys.\ Rev.\ D \textbf{98}, 056004 (2018), arXiv:1801.03917.
\\\textbf{[来源:} \texttt{docs/Factorization Theorem Relating Euclidean and Light-Cone Parton Distributions.pdf}\textbf{]}
——准TMD-PDF与光锥TMD-PDF之间的因子化定理的严格证明，确立软函数和CS核在LaMET中的角色。

\bibitem{Chu2023soft}
M.-H.~Chu \textit{et al.} (LPC),
``Lattice calculation of the intrinsic soft function and the Collins-Soper kernel,''
JHEP \textbf{08}, 172 (2023), arXiv:2306.06488.
\\\textbf{[来源:} \texttt{docs/Lattice Calculation of the Intrinsic Soft Function and the Collins-Soper Kernel.pdf}\textbf{]}
——内禀软函数和CS核的计算（含真空Wilson圈$Z_E$在CLS和MILC系综上的测量）。

\bibitem{Chu2023TMDWF}
M.-H.~Chu \textit{et al.} (LPC),
``Transverse-momentum-dependent wave functions of pion from lattice QCD,''
Phys.\ Rev.\ D \textbf{108}, 094513 (2023), arXiv:2302.09961.
\\\textbf{[来源:} \texttt{docs/Transverse-Momentum-Dependent Wave Functions of Pion from Lattice QCD.pdf}\textbf{]}

\bibitem{LPC2020soft}
Q.-A.~Zhang \textit{et al.} (LPC),
``Lattice-QCD calculations of TMD soft function through large-momentum effective theory,''
Phys.\ Rev.\ Lett.\ \textbf{125}, 192001 (2020), arXiv:2005.14572.
\\\textbf{[来源:} \texttt{docs/Lattice-QCD Calculations of TMD Soft Function Through Large-Momentum Effective Theory.pdf}\textbf{]}

\bibitem{HYPsmearing}
A.~Hasenfratz and F.~Knechtli,
``Flavor symmetry and the static potential with hypercubic blocking,''
Phys.\ Rev.\ D \textbf{64}, 034504 (2001), arXiv:hep-lat/0103029.
——HYP涂抹（Hypercubic Smeared Link）的原始提出。

\bibitem{Chen2025gluon}
C.~Chen \textit{et al.} (CLQCD, LPC),
``Unpolarized gluon PDF of the nucleon from lattice QCD in the continuum limit,''
(2025), arXiv:2510.26425.
\\\textbf{[来源:} \texttt{docs/Unpolarized gluon PDF of the nucleon from lattice QCD in the continuum limit.pdf}\textbf{]}

\bibitem{NoteGluonPDFs}
``Note of gluon PDFs,'' 内部笔记。
\\\textbf{[来源:} \texttt{docs/Note of gluon PDFs.pdf}\textbf{]}和\texttt{文档/note\_of\_gluon\_PDFs.pdf}

\end{thebibliography}

\vspace{0.5cm}
\noindent\rule{\textwidth}{0.5pt}
\small
\textbf{交叉参考建议：}本报告应结合同一\texttt{补充/}目录下的以下文件阅读：
\begin{itemize}[nosep]
    \item \texttt{格点QCD中的TMD\_PDF.pdf}——TMD-PDF中staple Wilson线的详细讨论
    \item \texttt{格点QCD中的光锥PDF与quasi\_PDF.pdf}——共线PDF中直Wilson线的讨论
    \item \texttt{格点QCD中的大动量有效理论.pdf}——LaMET框架中Wilson线的系统引入
    \item \texttt{格点QCD中的重夸克有效理论.pdf}——HQET中Wilson线与静态重夸克传播子的联系
    \item \texttt{格点QCD中的场强张量.pdf}——Clover项的plaquette离散化（Wilson圈的微观构成）
    \item \texttt{格点QCD中的关联函数.pdf}——强子关联函数中的Wilson线收缩
    \item \texttt{格点QCD中的费米子方案.pdf}——不同费米子方案对Wilson线重整化的影响
    \item \texttt{格点QCD中的UV涨落.pdf}——紫外发散与Wilson线自能的联系
\end{itemize}

\end{document}
